% Zone „Das kennst du schon“ – Prozentrechnung, Kl. 7, schwach
\hypertarget{zone}{}
\einheitenkopf*{Das kennst du schon}
\setcounter{aufgabe}{0}

\begin{aufgabe}{Ich kann Anteile am Streifen ablesen und zeigen. Wie viel Prozent des Streifens sind gefärbt?}
\swa{\leerfeld[\%]}{\streifen{50}{$0\,\%$}{$100\,\%$}}
\swa{\leerfeld[\%]}{\streifen{30}{$0\,\%$}{$100\,\%$}}
\swa{\leerfeld[\%]}{\streifen[20]{25}{$0\,\%$}{$100\,\%$}}
\anw{Färbe den Anteil im Streifen.}
\swa{$60\,\%$}{\streifen{0}{$0\,\%$}{$100\,\%$}}
\swa{$15\,\%$}{\streifen[20]{0}{$0\,\%$}{$100\,\%$}}
\anw{Die gefärbten 10\,\% des Streifens wiegen 4\,kg.}
\swz[1]{Wie viel kg sind 30\,\%?}{\streifenfrage{30}{$0\,$kg}{}}
\swz[1]{Wie viel kg wiegt der ganze Streifen, also 100\,\%?}{\streifen{10}{$0\,$kg}{$?$}}
\end{aufgabe}

\begin{aufgabe}{Ich kann einen Anteil als Bruch schreiben und kürzen. Schreibe den gefärbten Anteil als Bruch und kürze so weit wie möglich.}
\swz{2 von 6 Kästchen}{\bruchrechteck[5]{2}{6}}
\swz{5 von 10 Kästchen}{\bruchrechteck[5]{5}{10}}
\swz{6 von 15 Kästchen}{\bruchrechteck[5]{6}{15}}
\anw{Färbe zuerst den Anteil im Rechteck. Schreibe ihn dann als Bruch.}
\swz{In einer Tüte sind 3 rote und 7 grüne Bonbons. Welcher Anteil der Bonbons ist rot?}{\bruchrechteck[5]{0}{10}}
\end{aufgabe}

\begin{aufgabe}{Ich kann einen Anteil als Bruch schreiben und kürzen – weiter. Erweitere den Bruch auf den Nenner 100.}
\swz[1]{$\frac{1}{2}$}{\bruchrechteck[5]{1}{2}}
\swz[1]{$\frac{7}{10}$}{\bruchrechteck[5]{7}{10}}
\swz[1]{$\frac{9}{25}$}{\bruchrechteck[7]{9}{25}}
\swz[1]{$\frac{11}{20}$}{\bruchrechteck[7]{11}{20}}
\swz{$\frac{6}{40}$}{\bruchrechteck[8]{6}{40}}
\end{aufgabe}

\begin{aufgabe}{Ich kann Brüche und Dezimalzahlen ineinander umwandeln. Schreibe als Dezimalzahl.}
\swz[1]{$\frac{1}{2}$}{\zahlenstrahl[xmin=0,xmax=1.05,xstep=0.1,karo=0.8]{0.5/}}
\swz[1]{$\frac{3}{10}$}{\zahlenstrahl[xmin=0,xmax=1.05,xstep=0.1,karo=0.8]{0.3/}}
\swz[1]{$\frac{4}{5}$}{\zahlenstrahl[xmin=0,xmax=1.05,xstep=0.1,karo=0.8]{0.8/}}
\anw{Schreibe als Bruch.}
\swz[1]{$0{,}9$}{\zahlenstrahl[xmin=0,xmax=1.05,xstep=0.1,karo=0.8]{0.9/}}
\swz[1]{$0{,}09$}{\zahlenstrahl[xmin=0,xmax=0.105,xstep=0.01,karo=0.8]{0.09/}}
\end{aufgabe}

\begin{aufgabe}{Ich kann einen Anteil in Prozent angeben. Schreibe in Prozent.}
\swz[1]{$\frac{45}{100}$}{\streifen[20]{45}{$0\,\%$}{$100\,\%$}}
\swz[1]{$\frac{6}{10}$}{\streifen{60}{$0\,\%$}{$100\,\%$}}
\swz[1]{$0{,}65$}{\streifen[20]{65}{$0\,\%$}{$100\,\%$}}
\swz[1]{$0{,}07$}{\streifen[0]{7}{$0\,\%$}{$100\,\%$}}
\swz[1]{40\,\% der Kinder kommen zu Fuß, alle anderen fahren. Wie viel Prozent fahren?}{\streifen{40}{$0\,\%$}{$100\,\%$}}
\end{aufgabe}

\begin{aufgabe}{Ich kann auf eine Stelle nach dem Komma runden. Runde auf eine Stelle nach dem Komma.}
\swz[1]{$4{,}32$}{\zahlenstrahl[xmin=4,xmax=5.05,xstep=0.1,karo=0.8]{4.32/}}
\swz[1]{$7{,}68$}{\zahlenstrahl[xmin=7,xmax=8.05,xstep=0.1,karo=0.8]{7.68/}}
\swz[1]{$0{,}452$}{\zahlenstrahl[xmin=0,xmax=1.05,xstep=0.1,karo=0.8]{0.452/}}
\swz[1]{$19{,}96$}{\zahlenstrahl[xmin=19,xmax=20.05,xstep=0.1,karo=0.8]{19.96/}}
\anw{Der Taschenrechner zeigt diese Zahl. Runde sie.}
\swz[1]{$16{,}66666667$}{\zahlenstrahl[xmin=16,xmax=17.05,xstep=0.1,karo=0.8]{16.6667/}}
\end{aufgabe}

\begin{aufgabe}{Ich kann durch 100 teilen und mit Kommazahlen malnehmen. Teile durch 100.}
\begin{gleichungsraster}[1]
\gl{600 : 100} & \gl{85 : 100} \\
\gl{4 : 100} & \\
\end{gleichungsraster}\rasterende
\anw{Rechne.}
\begin{gleichungsraster}[1]
\gl{0{,}5 \cdot 18} & \gl{0{,}25 \cdot 12} \\
\gl{0{,}3 \cdot 40} & \\
\end{gleichungsraster}\rasterende
\end{aufgabe}

\begin{aufgabe}{Ich kann im Dreisatz erst auf 1 und dann hochrechnen. Fülle das Schema aus.}
\par\noindent\begin{minipage}[t]{0.48\linewidth}\vspace{0pt}
\begin{teile}\teil 5\,kg Äpfel kosten 10\,€. Wie viel kosten 3\,kg?\end{teile}
\begin{dreisatz}{Menge}{Preis}
\dsz{5\,\text{kg}}{10\,€}
\dsp{:5}{:5}
\dsz{1\,\text{kg}}{\dsleer}
\dsp{\cdot 3}{\cdot 3}
\dsz{3\,\text{kg}}{\dsleer}
\end{dreisatz}
\end{minipage}\hfill
\begin{minipage}[t]{0.48\linewidth}\vspace{0pt}
\begin{teile}\teil 3 Karten kosten 12\,€. Wie viel kosten 7 Karten?\end{teile}
\begin{dreisatz}{Anzahl}{Preis}
\dsz{3}{12\,€}
\dsp{:3}{:3}
\dsz{1}{\dsleer}
\dsp{\cdot 7}{\cdot 7}
\dsz{7}{\dsleer}
\end{dreisatz}
\end{minipage}\par\medskip
\noindent\begin{minipage}[t]{0.48\linewidth}\vspace{0pt}
\begin{teile}\teil 10\,m Stoff kosten 25\,€. Wie viel kosten 4\,m?\end{teile}
\begin{dreisatz}{Länge}{Preis}
\dsz{10\,\text{m}}{25\,€}
\dsp{:10}{:10}
\dsz{1\,\text{m}}{\dsleer}
\dsp{\cdot 4}{\cdot 4}
\dsz{4\,\text{m}}{\dsleer}
\end{dreisatz}
\end{minipage}\hfill
\begin{minipage}[t]{0.48\linewidth}\vspace{0pt}
\begin{teile}\teil 6 Liter Saft kosten 9\,€. Wie viel kosten 4 Liter? Trage auch die Rechenschritte ein.\end{teile}
\begin{dreisatz}{Menge}{Preis}
\dsz{6\,\text{l}}{9\,€}
\dsp{\phantom{:0}}{\phantom{:0}}
\dsz{\dsleer}{\dsleer}
\dsp{\phantom{:0}}{\phantom{:0}}
\dsz{4\,\text{l}}{\dsleer}
\end{dreisatz}
\end{minipage}\par
\end{aufgabe}

\begin{aufgabe}{Ich kann einen Bruchteil von einer Größe ausrechnen. Rechne erst durch den Nenner, dann mal den Zähler.}
\swz{$\frac{1}{2}$ von 30\,€}{\bruchrechteck[5]{1}{2}}
\swz{$\frac{1}{4}$ von 16\,kg}{\bruchrechteck[5]{1}{4}}
\swz{$\frac{1}{10}$ von 4,50\,€}{\bruchrechteck[5]{1}{10}}
\swz{$\frac{3}{5}$ von 20\,m}{\bruchrechteck[5]{3}{5}}
\end{aufgabe}

\begin{aufgabe}{Ich finde den Fehler bei einem Anteil. In der Klasse 7c sind 12 Mädchen und 13 Jungen. Emil soll angeben, welcher Anteil der Kinder Mädchen sind. Er schreibt:}
\rechnung{\text{Anteil der Mädchen} &= \tfrac{12}{13}}
\swz{Finde den Fehler, benenne ihn und schreibe den richtigen Anteil auf.}{\bruchrechteck[8]{12}{25}}
\swz{In einem Korb liegen 7 Äpfel und 13 Birnen. Welcher Anteil der Früchte sind Äpfel?}{\bruchrechteck[8]{0}{20}}
\end{aufgabe}
